\section*{第 \ref{ch:b1:mirrors-thin-lenses} 章——镜子与薄透镜}

\begin{solution}{exo:b1:mirrors-thin-lenses:1}
$\overline{SF} = \qty{-20}{cm}$；$1/\overline{SA'} = 1/\overline{SF} -
1/\overline{SA}$。当 \qty{60}{cm} 时：$\overline{SA'} = \qty{-30}{cm}$，
实像，$\gamma = -0.5$。当 \qty{30}{cm} 时：$\overline{SA'} = \qty{-60}{cm}$，
实像，$\gamma = -2$。当 \qty{10}{cm} 时：$1/\overline{SA'} = -1/20 + 1/10
= 1/20$，$\overline{SA'} = \qty{+20}{cm}$，位于镜后，虚像，
$\gamma = +2$（正立、放大）。
\end{solution}

\begin{solution}{exo:b1:mirrors-thin-lenses:2}
$\overline{OA} = \qty{-20}{cm}$：$1/\overline{OA'} = 1/5 - 1/20 = 3/20$，
$\overline{OA'} = \qty{+6.7}{cm}$，实像、倒立，$\gamma = -1/3$。
$\overline{OA} = \qty{-3.0}{cm}$：$1/\overline{OA'} = 1/5 - 1/3 = -2/15$，
$\overline{OA'} = \qty{-7.5}{cm}$，虚像、正立，$\gamma = +2.5$。
\end{solution}

\begin{solution}{exo:b1:mirrors-thin-lenses:3}
$V = 1/0.25 = \qty{+4.0}{\delta}$；$f' = 1/(-2.0) = \qty{-50}{cm}$；
$V = 3.0 - 1.0 = \qty{+2.0}{\delta}$，$f' = \qty{50}{cm}$。
\end{solution}

\begin{solution}{exo:b1:mirrors-thin-lenses:4}
$\overline{FA} = \qty{-15}{cm}$：$\overline{F'A'} = -f'^2/\overline{FA}
= +100/15 = \qty{6.7}{cm}$，位于 $F'$ 之后，即 $\overline{OA'} = \qty{16.7}{cm}$；
$\gamma = f'/\overline{FA} = -0.67$。按笛卡儿符号约定：$\overline{OA} =
\qty{-25}{cm}$，$1/\overline{OA'} = 1/10 - 1/25 = 3/50$，
$\overline{OA'} = \qty{16.7}{cm}$，$\gamma = 16.7/(-25) = -0.67$。
\end{solution}

\begin{solution}{exo:b1:mirrors-thin-lenses:5}
$\overline{SF} = \qty{+1.0}{m}$，$\overline{SA} = \qty{-20}{m}$：
$1/\overline{SA'} = 1 + 1/20$，$\overline{SA'} = \qty{0.95}{m}$，位于镜后
（虚像、正立）；$\gamma = -0.95/(-20) = +0.048$，像高
\qty{7.1}{cm}。大脑依据视大小判断距离：一辆
车在 \qty{0.95}{m} 处缩小二十倍，会被误读为
远处的车——远于 \qty{20}{m}。
\end{solution}

\begin{solution}{exo:b1:mirrors-thin-lenses:6}
$1/\overline{OA'} = 1/50 - 1/2000 = \qty{0.0195}{mm^{-1}}$，
$\overline{OA'} = \qty{51.3}{mm}$：透镜应向远离
感光元件的方向移动 \qty{1.3}{mm}。$\gamma = 51.3/(-2000) = -0.026$。若要 $\abs\gamma = 24/1800
= 1/75$，则 $\overline{OA} = f'(1/\gamma - 1) = 50 \times (-76) =
\qty{-3.8}{m}$。
\end{solution}

\begin{solution}{exo:b1:mirrors-thin-lenses:7}
$\gamma = -1.8/0.036 = -50$；$\overline{OA'} = \qty{5.0}{m}$，
$\overline{OA} = \overline{OA'}/\gamma = \qty{-10}{cm}$；
$1/f' = 1/5.0 + 1/0.10 = \qty{10.2}{m^{-1}}$，$f' = \qty{98}{mm}$。
像倒立（$\gamma < 0$），上下、左右方向都颠倒，幻灯片
应旋转 $180^\circ$ 放入。
\end{solution}

\begin{solution}{exo:b1:mirrors-thin-lenses:8}
物位于焦平面 $\Rightarrow$ 每个点发出的光线经透镜后
变为平行光；平面镜将其平行地反射回去（沿对称
方向）；透镜把平行光束会聚在焦平面上，位于
与物点关于光轴对称的点上：像位于物平面内，
倒立，$\gamma = -1$。平面镜倾斜 $\alpha$ 角时，
返回光束转过 $2\alpha$，像横向平移
$2\alpha f'$，但仍在同一平面内清晰成像。
\end{solution}

\begin{solution}{exo:b1:mirrors-thin-lenses:9}
$f' = (100^2 - 40^2)/(4 \times 100) = \qty{21}{cm}$。物像间最短距离
为 $4f' = \qty{84}{cm}$，此时 $\gamma = -1$。
\end{solution}

\begin{solution}{exo:b1:mirrors-thin-lenses:10}
$L_1$ 将无穷远处成像于 $F_1'$，位于 $L_1$ 之后 \qty{20}{cm}，即
位于 $L_2$ 之后 \qty{5}{cm}：虚物，$\overline{O_2A_1} = \qty{+5}{cm}$。
$1/\overline{O_2A'} = -1/10 + 1/5 = 1/10$：最终像位于
$L_2$ 之后 \qty{10}{cm}。中间像高 $h_1 = f_1'\alpha = 20\alpha$（cm）；
$\gamma_2 = 10/5 = 2$，故 $h = 40\alpha$：等效焦距
为 \qty{40}{cm}，镜筒长 $15 + 10 = \qty{25}{cm}$——这是一种长焦镜头。
\end{solution}

\begin{solution}{exo:b1:mirrors-thin-lenses:11}
$\gamma = -3 = -\overline{SA'}/\overline{SA}$，得 $\overline{SA'} =
3\overline{SA}$；两者均为实的，位于镜前：$\overline{SA} = -x$，
$\overline{SA'} = -3x$，且屏与灯的距离 $2x = \qty{80}{cm}$：
$x = \qty{40}{cm}$。镜子距灯 \qty{40}{cm}，
距屏 \qty{120}{cm}；$1/\overline{SF} = -1/40 - 1/120 = -1/30$：
$R = 2 \times 30 = \qty{60}{cm}$。
\end{solution}

\begin{solution}{exo:b1:mirrors-thin-lenses:12}
$\theta = 0.53^\circ = \qty{9.3e-3}{rad}$；像的直径 $f'\theta =
\qty{4.6}{mm}$。会聚功率 $10^3 \times \pi(0.05)^2 = \qty{7.9}{W}$；
像的面积 $\pi(2.3\times10^{-3})^2 = \qty{1.7e-5}{m^2}$；辐照度
$\qty{4.7e5}{W/m^2}$，为直射阳光的 $470$ 倍：纸张炭化并
燃烧——这就是取火镜。一般地，增益为
$(D/f'\theta)^2$。
\end{solution}

\begin{solution}{pb:b1:mirrors-thin-lenses:1}
\textbf{1.} $1/\overline{OA'} - 1/\overline{OA} = 1/f'$，
$\gamma = \overline{OA'}/\overline{OA}$；光轴沿光的传播方向取向，
距离自 $O$ 起算并带符号。

\textbf{2.} $p = -\overline{OA}$，$\overline{OA'} = D - p$：
$1/(D - p) + 1/p = 1/f_1'$，即 $p^2 - Dp + Df_1' = 0$；实根当且仅当
$D^2 - 4Df_1' \geq 0$，即 $D \geq 4f_1'$。

\textbf{3.} 两根之和为 $D$：$p_1 + p_2 = D$，故 $p_2 = D - p_1$——
物距与像距互换，两位置关于
中点对称。$d = p_2 - p_1 = \sqrt{D^2 - 4Df_1'}$，因此 $f_1' = (D^2 -
d^2)/4D$。

\textbf{4.} $f_1' = (6400 - 1600)/320 = \qty{15.0}{cm}$.

\textbf{5.} 只有两次清晰成像之间镜筒的位移 $d$
与物屏距离 $D$ 参与计算：二者都不需要知道
$O$ 位于镜筒内何处。

\textbf{6.} $\partial f_1'/\partial D = 1/4 + d^2/4D^2 = 0.3125$,
$\partial f_1'/\partial d = -d/2D = -0.25$；于是
$u(f_1')^2 = (0.3125 \times 0.2)^2 + (0.25 \times 0.5)^2 =
\qty{0.0195}{cm^2}$，且 $u(f_1') = \qty{0.14}{cm}$：
$f_1' = \qty{15.00\pm0.14}{cm}$.

\textbf{7.} $p = (80 \mp 40)/2 = 20$ 或 \qty{60}{cm}：$\gamma =
-60/20 = -3$（幻灯片靠近透镜，像被放大：投影仪）和
$\gamma = -1/3$.

\textbf{8.} $D_{\min} = 4f_1' = \qty{60.0}{cm}$, $\gamma = -1$.

\textbf{9.} $\gamma = -1.80/0.036 = -50$.

\textbf{10.} $\overline{OA'} = f_1'(1 - \gamma) = 15 \times 51 =
\qty{765}{cm} = \qty{7.65}{m}$; $\overline{OA} = \overline{OA'}/\gamma =
\qty{-15.3}{cm}$.

\textbf{11.} $\gamma < 0$：实像被旋转 $180^\circ$
绕光轴，故幻灯片应旋转 $180^\circ$ 后插入（上下颠倒
且左右镜像）。

\textbf{12.} $\overline{OA'} = \qty{1000}{cm}$：$1/\overline{OA} = 1/1000
- 1/15$，$\overline{OA} = \qty{-15.23}{cm}$：透镜移动 \qty{0.7}{mm}
朝向幻灯片。$\gamma = 1000/(-15.23) = -65.7$：
画面宽 \qty{2.36}{m}。

\textbf{13.} 对笛卡儿公式求微分：$-\dd\overline{OA'}/\overline{OA'}^2
+ \dd\overline{OA}/\overline{OA}^2 = 0$，故 $\dd\overline{OA'} =
(\overline{OA'}/\overline{OA})^2\,\dd\overline{OA} = \gamma^2\,\dd\overline{OA}$.

\textbf{14.} $2500 \times \qty{0.10}{mm} = \qty{25}{cm}$：像
随胶片受热而离焦——因此需要用玻璃框固定的幻灯片，或
能跟随胶片移动的自动对焦装置。

\textbf{15.} 由光阑至像点作相似三角形：
弥散斑 $= 30 \times 0.20/7.65 = \qty{0.78}{mm}$；自 \qty{5}{m} 处看，其张角
为 $\qty{1.6e-4}{rad} < \qty{3e-4}{rad}$，不易察觉。

\textbf{16.} 幻灯片面积 $8.64\times10^{-4}\ \unit{m^2}$：$E_{\text{幻灯片}}
= 2.0/8.64\times10^{-4} = \qty{2.3e3}{W/m^2}$；屏幕面积为幻灯片面积的 $\gamma^2$
倍，即 $\qty{2.16}{m^2}$：$E_{\text{屏幕}} = \qty{0.93}{W/m^2}$；
功率相同，面积之比为 $\gamma^2 = 2500$。

\textbf{17.} $\overline{O_2A'} = \qty{1200}{cm}$：$1/\overline{O_2A} =
1/1200 + 1/30 = 0.03417$，$\overline{O_2A} = \qty{+29.3}{cm}$：
中间像位于 $L_2$ \emph{之后} \qty{29.3}{cm} 处——对
$L_2$ 而言为虚物（来自 $L_1$ 的光线朝其会聚）。

\textbf{18.} $\overline{O_1A_1'} = 8.0 + 29.3 = \qty{37.3}{cm}$;
$1/\overline{O_1A} = 1/37.3 - 1/15 = -0.0398$, $\overline{O_1A} =
\qty{-25.1}{cm}$.

\textbf{19.} $\gamma_1 = 37.3/(-25.1) = -1.49$; $\gamma_2 = 1200/29.3 =
41.0$; $\gamma = -60.9$.

\textbf{20.} 宽度 $60.9 \times 36 = \qty{2.19}{m}$。仅用 $L_1$ 时：
$\gamma = 1 - \overline{OA'}/f_1' = 1 - 1208/15 = -79.5$，宽度
\qty{2.86}{m}.

\textbf{21.} 以 $\alpha$ 入射的光束：$L_1$ 将其会聚于焦平面上高度 $h_1 = f_1'\alpha$
处，$\overline{O_2A_1} = f_1' - e$ 即该焦平面与 $L_2$ 的距离；
$L_2$ 将该点成像于 $\overline{O_2A'}$ 处，其中 $1/\overline{O_2A'}
= 1/f_2' + 1/(f_1' - e)$，放大率为 $\gamma_2 = \overline{O_2A'}/(f_1'
- e)$，故 $h = \gamma_2 f_1'\alpha$，且 $f_{\mathrm{eq}} = \gamma_2 f_1'
= \dfrac{f_1' f_2'}{f_1' + f_2' - e}$，其倒数即所给出的
公式。$1/f_{\mathrm{eq}} = 1/15 - 1/30 + 8/450 = \qty{0.0511}{cm^{-1}}$，
$f_{\mathrm{eq}} = \qty{19.6}{cm}$.

\textbf{22.} $2.19/2.86 = 0.77$ 与 $15/19.6 = 0.77$。仅为近似，
因为幻灯片位于有限距离处，且两个投射距离
分别自不同的透镜量起；若物体位于无穷远，该比值
将是精确的。

\textbf{23.} 位于凹透镜之后、$p > 0$ 处的虚物：
$1/\overline{O_2A'} = 1/p + 1/f_2' = 1/p - 1/\abs{f_2'}$，为正（实
像）当且仅当 $p < \abs{f_2'}$。

\textbf{24.} 功率相同，画面面积为原来的 $(2.19/2.86)^2 = 0.59$
倍，辐照度 $\times 1.7$。

\textbf{25.} 转换器使焦距乘以
$f_{\mathrm{eq}}/f_1' = 1.3$；在给定投射距离下，画面
各方向上缩小 $1.3$ 倍，亮度增加 $1.7$ 倍——是一种
远投镜头。
\end{solution}
